\documentclass[12pt,a4paper]{article}
\usepackage[utf8]{inputenc}
\usepackage[T2A]{fontenc}
\usepackage[russian]{babel}
\usepackage{amsmath,amssymb,amsthm}
\usepackage{booktabs,array}
\usepackage{geometry}
\usepackage{microtype}
\geometry{margin=2.3cm}

\newtheorem{proposition}{Предложение}
\newtheorem{nogocriterion}{Критерий провала}
\DeclareMathOperator{\Spec}{Spec}

\title{Spin-menu Dirac и determinant-line gate:\\
последняя спектральная попытка вывести family-selector}
\author{}
\date{5 августа 2026 г.}

\begin{document}
\maketitle

\section{Внешний математический контроль}

Теория determinant line для семейств операторов Дирака связывает
экспоненцированную $\eta$-инварианту с элементом determinant line, а её
глобальное изменение --- с голономией естественной связи. Зависимость
$\eta$-инварианты от spin-структуры также является реальным спектральным
эффектом. Поэтому остаётся законный вопрос: может ли семейство четырёх Dirac-
операторов на spin-меню породить недостающие off-diagonal family-переходы.

Однако отдельные determinant- или $\eta$-значения прикреплены к различным
spin-структурам как скаляры/фазы. Без дополнительного parallel transport они
задают только диагональный оператор. Чтобы получить incidence matrix, необходимо
ввести связь между determinant fibers.

\section{Минимальный cellular Dirac}

Четыре spin-состояния образуют квадрат $F_2^2$. Его две группы рёбер
соответствуют переключению $\mathbb{RP}^3$- и $S^1$-spin-структур. Для
ориентированной incidence matrix $B$ естественный дискретный оператор равен
\begin{equation}
 D_{cell}=\begin{pmatrix}0&B\\ B^*&0\end{pmatrix},
 \qquad \Delta_0=BB^*.
\end{equation}
При весах $w_3=\pi^{-1}$ и $w_1=(2\pi)^{-1}$ получаем
\begin{equation}
 \Spec(\Delta_0)=
 \left\{0,\frac1\pi,\frac2\pi,\frac3\pi\right\}.
\end{equation}
Оператор точно коммутирует с обеими факторными трансляциями. Следовательно,
cellular Dirac square возвращает уже проверенный факторный лапласиан: он
различает семейные метки, но не создаёт mixing.

\section{Связь на determinant line}

После vertex rephasing физической характеристикой связи на квадрате остаётся
голономия вокруг plaquette. Для projective representation группы
$\mathbb Z_2^2$ существуют два топологических класса:
\begin{equation}
 \operatorname{Hol}_{\square}=+1
 \quad\text{или}\quad
 \operatorname{Hol}_{\square}=-1.
\end{equation}
Первый класс gauge-эквивалентен тривиальной commuting connection. Второй даёт
\begin{equation}
 UV=-VU
\end{equation}
и совпадает с ранее исследованным $\pi$-flux.

В явном weighted cellular расчёте $\pi$-flux даёт две собственные величины
\begin{equation}
 0.1215835576,\qquad 0.8333461010,
\end{equation}
каждую с кратностью два. Кроме того, magnetic Laplacian смешивает uniform-
синглет и триплет. Поэтому determinant-line transport не создаёт нового
третьего маршрута: он возвращает либо commuting case, либо even-multiplicity
projective case.

\begin{proposition}[Spectral-selector dichotomy]
Для минимального квадратного spin-меню canonical Dirac/determinant construction
имеет только три формы чтения: диагональные determinant/$\eta$-веса,
тривиальная commuting connection или нетривиальный $\pi$-flux. Ни одна из них
не выбирает outside-$D_8$ incidence-оператор, порождающий полную $M_3$-алгебру.
\end{proposition}

\section{Финальный семейный gate текущего носителя}

Spectral flow способен присвоить ориентированному пути целое число или фазу,
но сам граф путей и его edge magnitudes являются дополнительными данными.
После факторизации vertex gauge остаётся только уже учтённая plaquette-
голономия. Следовательно, ни cellular Dirac, ни $\eta$-инварианты, ни
determinant-line holonomy не выводят искомый selector из текущего носителя.

\begin{nogocriterion}[Закрытие intrinsic-selector ветви]
В пределах данных $\mathbb{RP}^3\times S^1$, четырёх spin-структур, их
cohomology ring, product metric и canonical determinant line недостающий
full-$M_3$ incidence-оператор не выводится. Переоткрытие требует новой
неплоской boundary connection, заранее фиксирующей граф, edge magnitudes и
различное назначение up/down-секторам.
\end{nogocriterion}

Этот результат закрывает не все возможные flavour-теории, а intrinsic-
selector программу текущего конечного носителя. Двенадцать algebraically
рабочих операторов существуют, но ни один не выделяется имеющимися
геометрическими и спектральными данными.

Машинный протокол сохранён в
\texttt{s2t\_spin\_menu\_dirac\_determinant\_line\_audit.py} и
\texttt{s2t\_spin\_menu\_dirac\_determinant\_line\_results.json}.

\end{document}